\chapter{Business Intelligence (BI)}
\label{cap:bi}

\begin{edital}
Conceitos, fundamentos, técnicas e métodos de BI; sistemas de suporte à
decisão (SSD) e gestão de conteúdo; arquitetura e aplicações de data
warehouse com ETL e OLAP; data warehouse e data mining; visualização (BD
individuais e cubos); mapeamento das fontes de dados; arquitetura de BI.
\end{edital}
\peso{MÉDIO --- costuma cair junto com Banco de Dados.}

\section{O que é BI}

Conjunto de conceitos e ferramentas para transformar \textbf{dado} em
\textbf{informação} e apoiar a \textbf{tomada de decisão}. Pirâmide clássica:
dado $\to$ informação $\to$ conhecimento $\to$ inteligência/sabedoria.

A pergunta que fundou o assunto é mais concreta: \emph{por que não rodar o
relatório direto no banco do sistema?} Por três motivos que aparecem, um a um,
nos enunciados de cenário. Primeiro, \textbf{concorrência}: uma consulta que
varre cinco anos de vendas segura recursos que o sistema transacional precisa
para atender o caixa \emph{agora}. Segundo, \textbf{história}: o banco
transacional guarda o estado \emph{atual} --- ele sabe onde o cliente mora, não
onde morava quando comprou. Terceiro, e o mais caro, \textbf{integração}: os
dados estão em vários sistemas que discordam entre si (a mesma unidade com dois
códigos, a mesma data em dois formatos). O BI existe para resolver os três de
uma vez, e é isso que o resto do capítulo descreve.

\begin{conceito}[A arquitetura de BI --- o caminho do dado, ponta a ponta]
\begin{center}
fontes (OLTP, planilhas, externas) $\to$ \textbf{extração} $\to$
\emph{staging} $\to$ \textbf{transformação} $\to$ \textbf{DW} (e seus
\emph{data marts}) $\to$ \textbf{OLAP/cubo} $\to$ visualização
\end{center}

Cada etapa existe por um motivo, e é o motivo que a banca cobra:

\begin{itemize}
\item A \textbf{extração} lê das fontes sem incomodá-las mais do que o
necessário (janela de carga, leitura incremental do que mudou).
\item A \emph{\textbf{staging area}} é o pouso intermediário: dado bruto já
fora da fonte, ainda não tratado. Existe para que a fonte seja liberada
depressa e para que o tratamento possa ser refeito sem voltar a ela.
\item A \textbf{transformação} é onde a integração acontece --- é a etapa cara,
não a cópia.
\item O \textbf{DW} guarda o resultado: integrado, histórico e
\textbf{não volátil} (não se atualiza um fato passado; acrescenta-se).
\item \textbf{OLAP e visualização} são \textbf{consumo}. Nenhum dos dois
transforma dado --- quem transforma é o ETL, e essa separação é a resposta de
vários itens.
\end{itemize}

O \textbf{data mart} é o recorte do DW por assunto ou departamento (vendas,
RH). Ele não é uma cópia menor por economia: é a entrega no vocabulário de quem
vai usar.
\end{conceito}

\section{Data Warehouse (DW)}

Repositório central de dados \textbf{integrados, históricos e orientados a
assunto}, para análise (não para transação). Características de Inmon:
\textbf{orientado a assunto, integrado, não volátil, variante no tempo}.

\begin{itemize}
\item \textbf{Data Mart:} recorte departamental do DW.
\item \textbf{DW $\times$ Data Lake:} DW guarda dado \textbf{tratado e
modelado} (schema-on-write); Lake guarda dado \textbf{bruto} estruturado e
não estruturado (schema-on-read).
\item \textbf{Staging area:} área intermediária onde o ETL trata os dados.
\item \textbf{Lakehouse:} une a governança do DW com a flexibilidade do
Lake. Camadas típicas: \textbf{Bronze} (bruto) $\to$ \textbf{Silver} (limpo)
$\to$ \textbf{Gold} (dimensional, pronto para BI). Machine Learning consome o
bruto; BI consome o Gold.
\end{itemize}

\begin{pegadinha}
Não diga que o Data Lake ``substitui'' o DW, nem que ele guarda só dado
estruturado. Lakehouse é o que combina os dois --- não é sinônimo de nenhum
dos dois isoladamente.
\end{pegadinha}

\subsection{Inmon $\times$ Kimball --- por onde se começa}

As duas escolas discordam sobre \textbf{qual peça se constrói primeiro}, e a
discordância tem consequência prática.

\begin{conceito}[Top-down $\times$ bottom-up]
\textbf{Inmon (top-down).} Constrói-se primeiro o \textbf{DW corporativo},
normalizado (até a 3FN), como fonte única de verdade da organização; os
\textbf{data marts} dimensionais são \emph{derivados} dele. Integra melhor e
resiste melhor à mudança, mas o primeiro relatório demora --- a área usuária
espera o modelo corporativo ficar de pé.

\textbf{Kimball (bottom-up).} Começa-se por um \textbf{data mart dimensional}
de um processo de negócio (vendas, por exemplo), já em esquema estrela, e vão
se acrescentando outros. Entrega valor cedo. A integração não vem do modelo
central: vem das \textbf{dimensões conformadas} --- a mesma dimensão
\texttt{Produto}, com a mesma chave e o mesmo significado, compartilhada por
todos os marts (a \emph{bus architecture}).

\textbf{O risco de cada um é o espelho do outro.} Em Inmon, o projeto longo que
nunca chega ao usuário; em Kimball, o \textbf{silo}: marts que não conversam
porque as dimensões não foram conformadas. Por isso ``dimensão conformada''
não é detalhe de vocabulário --- é a peça que faz o bottom-up funcionar.
\end{conceito}

\section{ETL --- a ponte para o DW}

\textbf{Extract $\to$ Transform $\to$ Load.} Objetivo principal: extrair de
várias fontes, transformar em formato padronizado e consistente, e carregar
no DW. \textbf{Não é} ``criar dashboards'' nem ``treinar modelo de ML'' ---
isso vem depois, como consumo do dado já carregado. ETL $\times$ ELT: ver
Capítulo \ref{cap:banco-dados} (a posição do T é a mesma pegadinha).

\begin{conceito}[O que acontece em cada fase --- e por que a FGV pergunta isso]
O item típico descreve uma tarefa e pergunta \textbf{a que fase ela pertence}.
Como quase tudo o que é interessante mora no \emph{T}, o critério tem de ser
mais fino que ``transformar é mudar o dado''.

\begin{center}
\begin{tabular}{@{}p{2.1cm}p{9.6cm}@{}}
\toprule
\textbf{Fase} & \textbf{O que é dela} \\
\midrule
\textbf{Extract} & ler das fontes (banco, arquivo, API), tipicamente só o que
mudou desde a última carga. \textbf{Sem regra de negócio} --- ela lê, não julga \\
\textbf{Transform} & \textbf{limpeza} (nulos, formatos, valores inválidos),
\textbf{padronização/conformidade} (o de-para que faz dois sistemas falarem a
mesma língua), \textbf{deduplicação}, \textbf{junção} de fontes diferentes,
\textbf{derivação} de medidas calculadas, geração da \textbf{chave
\emph{surrogate}} e aplicação da regra de SCD \\
\textbf{Load} & gravar no destino, \textbf{dimensões antes da fato} (a fato
referencia as chaves das dimensões), carga completa $\times$ incremental,
índices e agregados no fim \\
\bottomrule
\end{tabular}
\end{center}

\textbf{A regra prática:} se a tarefa \emph{decide} alguma coisa sobre o
conteúdo --- padronizar, deduplicar, casar duas fontes, calcular ---, é
\textbf{Transform}. Ler é Extract; gravar é Load.

\textbf{ELT} inverte a ordem por um motivo econômico: quando o destino tem
poder de processamento de sobra (nuvem, \emph{lakehouse}), carrega-se o bruto e
transforma-se \textbf{dentro} dele, com SQL. Não é uma versão moderna do ETL
que torne a outra obsoleta --- é a mesma decisão de onde gastar o
processamento.
\end{conceito}

\section{OLAP e modelagem dimensional}

OLAP (analítico) $\times$ OLTP (transacional): ver Capítulo
\ref{cap:banco-dados}.

\textbf{Cubo:} dados organizados em \textbf{dimensões} (tempo, produto,
local) e \textbf{métricas/fatos} (vendas, quantidade).

\begin{conceito}[Modelo relacional $\times$ modelo multidimensional]
São \textbf{duas formas de organizar o mesmo dado}, para finalidades
diferentes --- e a FGV cobra a comparação diretamente.

\begin{center}
\begin{tabular}{@{}p{2.6cm}p{4.2cm}p{4.6cm}@{}}
\toprule
& \textbf{Relacional (OLTP)} & \textbf{Multidimensional (OLAP)} \\
\midrule
Organiza em & tabelas normalizadas (até 3FN) & fatos $\times$ dimensões (cubo) \\
Otimizado para & \textbf{escrita} curta e concorrente, sem redundância & \textbf{leitura} analítica, agregação em massa \\
Redundância & evitada & \textbf{aceita de propósito} (desnormalização) \\
Pergunta típica & ``qual o saldo do cliente X agora?'' & ``como as vendas evoluíram por região e trimestre?'' \\
\bottomrule
\end{tabular}
\end{center}

Formas de \textbf{armazenar} o cubo: \textbf{ROLAP} (o dado fica em tabelas
relacionais e o cubo é montado por consulta --- escala melhor),
\textbf{MOLAP} (cubo pré-calculado em estrutura própria --- consulta mais
rápida, carga mais pesada) e \textbf{HOLAP} (híbrido: detalhe no relacional,
agregados pré-calculados).
\end{conceito}

\begin{pegadinha}
Duas trocas frequentes: dizer que o modelo multidimensional serve para
\textbf{transação} (é o relacional/OLTP) e inverter ROLAP com MOLAP ---
lembre que o \textbf{M} de MOLAP é de \emph{multidimensional}, o cubo
pré-calculado. Absoluto a descartar: ``ferramentas de BI são incompatíveis
com tabelas normalizadas'' --- não são; a desnormalização é escolha de
desempenho, não exigência.
\end{pegadinha}

\begin{center}
\begin{tabular}{@{}p{3cm}p{8.6cm}@{}}
\toprule
\textbf{Operação OLAP} & \textbf{O que faz} \\
\midrule
Drill-down & desce ao detalhe (ano $\to$ mês $\to$ dia) \\
Roll-up (drill-up) & sobe à agregação \\
Slice & fatia \textbf{uma} dimensão (fixa um valor) \\
Dice & subcubo --- filtro em \textbf{várias} dimensões e valores \\
Pivot (rotate) & gira a visão \\
\bottomrule
\end{tabular}
\end{center}

\begin{conceito}[Star Schema $\times$ Snowflake Schema]
\textbf{Star Schema:} uma tabela fato + dimensões \textbf{desnormalizadas}.
Menos joins, mais rápido para BI --- é o padrão recomendado por Kimball.

\textbf{Snowflake Schema:} dimensões \textbf{normalizadas} em várias tabelas.
Mais joins, mais lento, mas economiza espaço e ajuda integridade.
\end{conceito}

\begin{pegadinha}
Star $\times$ Snowflake invertido: a FGV chama de ``estrela'' a dimensão
normalizada em várias tabelas (isso é floco de neve/Snowflake), ou diz que
Snowflake é o padrão recomendado por Kimball (é o Star). Também drill-down
$\leftrightarrow$ roll-up e OLAP $\leftrightarrow$ OLTP invertidos.
\end{pegadinha}

\textbf{Exemplo trabalhado --- montando o esquema estrela de vendas.} O projeto
dimensional tem uma ordem, e ela é sempre a mesma:

\begin{enumerate}
\item \textbf{Escolher o processo de negócio:} a venda no varejo.
\item \textbf{Declarar o grão:} ``\textbf{uma linha por item de um pedido}''.
É a decisão que governa todo o resto --- veja a seção seguinte.
\item \textbf{Escolher as dimensões:} tudo o que responde a ``por qual
recorte eu quero olhar?'' --- \texttt{Tempo}, \texttt{Produto}, \texttt{Loja},
\texttt{Cliente}.
\item \textbf{Escolher as medidas:} o que é \emph{numérico e somável} no grão
declarado --- \texttt{quantidade}, \texttt{valor\_unitário},
\texttt{valor\_total}, \texttt{desconto}.
\end{enumerate}

O resultado é uma \textbf{tabela fato} que contém \textbf{só duas coisas}: as
\textbf{chaves estrangeiras} para as dimensões e as \textbf{medidas
numéricas}. Nada de texto descritivo --- o nome do produto mora na dimensão
\texttt{Produto}, não na fato. A única exceção com nome próprio é a
\textbf{dimensão degenerada}: o \texttt{número\_do\_pedido} fica na fato,
porque é um identificador que não tem atributo nenhum para formar dimensão.

Agora a diferença que a banca mede em \emph{joins}. Para responder ``vendas do
\textbf{departamento} X no trimestre'':

\begin{center}
\begin{tabular}{@{}p{2.5cm}p{9.2cm}@{}}
\toprule
\textbf{Estrela} & \texttt{Produto} guarda categoria e departamento como
colunas dela mesma (desnormalizada): \textbf{2 joins} --- fato $\to$ Produto,
fato $\to$ Tempo \\
\textbf{Floco de neve} & \texttt{Produto} $\to$ \texttt{Categoria} $\to$
\texttt{Departamento}, cada uma em sua tabela: \textbf{4 joins} para a mesma
resposta \\
\bottomrule
\end{tabular}
\end{center}

É daí que sai tudo o mais: menos \emph{joins} $\to$ consulta mais rápida $\to$
esquema estrela como padrão do BI. O floco de neve troca velocidade por
economia de espaço e por integridade da hierarquia --- uma escolha legítima,
só não a padrão.

\subsection{Granularidade (o grão da fato)}

\textbf{Granularidade é o que uma linha da tabela fato representa.} Definir o
grão é a primeira decisão do projeto dimensional, e ela é irreversível na
prática: tudo o mais (dimensões, medidas, agregações) é escolhido depois
dela.

\begin{center}
\begin{tabular}{@{}p{3.4cm}p{7.6cm}@{}}
\toprule
\textbf{Grão} & \textbf{Consequência} \\
\midrule
\textbf{Fino} (detalhado) --- ``uma linha por atendimento'' & \textbf{amplia} as análises possíveis (dá para agregar depois em qualquer nível); volume maior \\
\textbf{Grosso} (agregado) --- ``uma linha por unidade por dia'' & tabela menor e consulta mais rápida, mas o detalhe está \textbf{perdido para sempre} --- não dá para desagregar \\
\bottomrule
\end{tabular}
\end{center}

Regra de Kimball: \textbf{declare o grão no nível mais atômico que a fonte
permitir}. Agregados são construídos \textbf{a partir} do detalhe, como
tabelas adicionais --- nunca no lugar dele.

\begin{pegadinha}
A FGV inverte a consequência: diz que escolher o grão mais detalhado
\textbf{restringe} as análises (é o contrário --- quem restringe é o
agregado) ou que a granularidade é definida pela quantidade de dimensões (é
definida pelo que \textbf{uma linha da fato representa}). Não confunda
\textbf{grão} com \textbf{nível de agregação de uma consulta}: o grão é do
modelo, a agregação é do drill-up.
\end{pegadinha}

\subsection{Tipos de fato e dimensão}

\begin{itemize}
\item \textbf{Fato aditivo:} soma em \textbf{todas} as dimensões (ex.:
vendas).
\item \textbf{Fato semiaditivo:} \textbf{não} soma em todas as dimensões,
tipicamente não soma no tempo (ex.: saldo, estoque).
\item \textbf{Fato não aditivo:} razão/percentual, nunca soma.
\item \textbf{Dimensão degenerada:} atributo de identificação (ex.: nº do
pedido) que fica na tabela fato, sem dimensão própria.
\end{itemize}

\begin{pegadinha}
A banca troca semiaditivo por aditivo --- lembre que semiaditivo tipicamente
\textbf{não} soma no eixo do tempo (saldo bancário de hoje não é a soma dos
saldos anteriores).
\end{pegadinha}

\subsection{Dimensões lentamente mutantes (SCD)}

Atributo de dimensão muda com o tempo (o cliente troca de cidade, o produto
troca de categoria). \textbf{SCD} (\textit{Slowly Changing Dimension}) é a
técnica que decide o que fazer com o histórico:

\begin{center}
\begin{tabular}{@{}p{1.9cm}p{4cm}p{5.2cm}@{}}
\toprule
\textbf{Tipo} & \textbf{O que faz} & \textbf{Efeito no histórico} \\
\midrule
\textbf{Tipo 1} & \textbf{sobrescreve} o valor antigo & histórico \textbf{perdido}; os fatos passados passam a ser lidos com o valor novo \\
\textbf{Tipo 2} & \textbf{insere um novo registro} versionado (nova chave \textit{surrogate}, com \texttt{data\_início}/\texttt{data\_fim} ou \textit{flag} de corrente) & \textbf{preserva o histórico completo} --- cada fato continua ligado à versão vigente na época \\
\textbf{Tipo 3} & guarda o \textbf{valor anterior em outra coluna} & preserva só \textbf{uma} mudança (o ``valor anterior''), não a série toda \\
\bottomrule
\end{tabular}
\end{center}

É o \textbf{tipo 2} que responde ao cenário clássico: ``as vendas antigas
devem continuar associadas à cidade em que o cliente morava na época''. E é
ele que justifica a \textbf{chave \textit{surrogate}} (substituta) da
dimensão: como a mesma chave natural passa a ter várias versões, a fato
precisa apontar para uma chave artificial que identifique a \textbf{versão},
não a entidade.

\begin{pegadinha}
Inversão de número é a pegadinha inteira aqui: oferecer o \textbf{tipo 1}
(sobrescrever) para um cenário que pede histórico, ou o \textbf{tipo 2} para
um cenário de correção de erro de digitação (aí o certo \emph{é} sobrescrever
--- tipo 1). O tipo 3 aparece como distrator ``quase certo'': preserva
histórico, mas só o \textbf{valor imediatamente anterior}.
\end{pegadinha}

\section{Sistemas de Suporte à Decisão (SSD/DSS)}

Apoiam gestores em problemas \textbf{estruturados, semiestruturados e não
estruturados} --- oferecem flexibilidade para vários contextos. Não são ``só
para problemas estruturados'' nem ``só não estruturados''.

\section{Mapeamento de fontes de dados (boas práticas)}

\begin{itemize}
\item \textbf{Fazer:} entrevistar stakeholders, analisar sistemas
existentes, avaliar \textbf{qualidade} e adequação do dado, documentar fontes
(governança). Referência: \textbf{DAMA-DMBOK}, dados mestres e de
referência.
\item \textbf{NÃO fazer:} coletar \textbf{tudo sem discriminação}, incluindo
dados inconsistentes/irrelevantes ``para maximizar quantidade'' --- essa é a
prática \textbf{não recomendada} que a FGV pede como resposta.
\end{itemize}

\section{Data Mining e visualização}

\textbf{Data Mining:} descobrir padrões, correlações e conhecimento não
trivial em grandes volumes de dados.

\begin{center}
\begin{tabular}{@{}p{3.2cm}p{3.2cm}p{4.4cm}@{}}
\toprule
\textbf{Tarefa} & \textbf{Tem rótulo?} & \textbf{Cenário típico} \\
\midrule
Classificação & \textbf{sim} (supervisionada) & prever se o cliente vai inadimplir (classes conhecidas) \\
Regressão & \textbf{sim} (supervisionada) & prever um \textbf{valor numérico} (faturamento do mês) \\
Clusterização & \textbf{não} (não supervisionada) & \textbf{descobrir grupos} de clientes com comportamento parecido, \emph{sem} rótulo prévio \\
Regras de associação & não & \textbf{o que é comprado junto} (cesta de compras) \\
\bottomrule
\end{tabular}
\end{center}

\textbf{Regras de associação:} \textbf{suporte} = frequência do conjunto no
total de transações; \textbf{confiança} = força da implicação A$\to$B. A FGV
confunde os dois.

\begin{pegadinha}
O par que a banca inverte é \textbf{classificação} $\times$
\textbf{clusterização}: se o enunciado diz ``\textbf{sem rótulos prévios}''
ou ``descobrir grupos'', é clusterização (não supervisionada); se as classes
já existem e o modelo aprende a atribuí-las, é classificação
(supervisionada). ``Produtos comprados juntos'' é \textbf{regra de
associação}, nunca clusterização.
\end{pegadinha}

\subsection{CRISP-DM --- o ciclo de um projeto de mineração}

\begin{center}
\begin{tabular}{@{}p{0.6cm}p{3.6cm}p{6.8cm}@{}}
\toprule
& \textbf{Fase} & \textbf{O que acontece} \\
\midrule
1 & Entendimento do \textbf{negócio} & objetivos e critérios de sucesso do negócio (é por aqui que começa, não pelos dados) \\
2 & Entendimento dos \textbf{dados} & coleta inicial, exploração, qualidade \\
3 & \textbf{Preparação} dos dados & limpeza, seleção, integração, formatação \\
4 & \textbf{Modelagem} & escolha da técnica, treino, ajuste de parâmetros \\
5 & \textbf{Avaliação} & os resultados atendem aos \textbf{objetivos de negócio}? \\
6 & \textbf{Implantação} & colocar o modelo em operação, monitorar \\
\bottomrule
\end{tabular}
\end{center}

O ciclo é \textbf{iterativo}, não uma cascata: volta-se de qualquer fase à
anterior, e a seta entre implantação e entendimento do negócio fecha o
círculo.

\begin{pegadinha}
Não confunda a \textbf{avaliação} (fase 5, contra os objetivos de
\textbf{negócio}) com a validação técnica do modelo, que é parte da
\textbf{modelagem} (fase 4). Depois da avaliação vem a \textbf{implantação}
--- se a alternativa disser que depois de avaliar volta-se sempre à
preparação, é extrapolação. Outro distrator: começar o ciclo pelo
entendimento \textbf{dos dados} --- começa pelo \textbf{negócio}.
\end{pegadinha}

\textbf{Dashboards e relatórios interativos:} camada de visualização (Power
BI --- com recurso de \textit{drillthrough} entre páginas ---, Qlik, etc.).

\begin{jacaiu}
\textbf{Em prova real da FGV:} SSD para os três tipos de problema e a prática
\emph{não} recomendada no mapeamento de fontes (coletar tudo sem
discriminar) --- \textbf{Dataprev 2024}; Star $\times$ Snowflake em cenário
de varejo/DW (contagem de \textit{joins}), o que a tabela fato contém
(chaves estrangeiras + medidas numéricas), Data Mart como recorte
departamental e o que pertence à fase de \textbf{Transformação} do ETL ---
\textbf{ALERO 2026}; arquitetura \textbf{Lakehouse} e \textbf{dice} (filtro
em várias dimensões) --- \textbf{TJ-RJ}; \textbf{suporte} $\times$
\textbf{confiança} em regras de associação, relacional $\times$
multidimensional e \textit{drillthrough} no Power BI --- \textbf{MPU}. O par
ETL $\times$ ELT caiu na \textbf{Dataprev 2024}, como item de Banco de Dados
(Capítulo \ref{cap:banco-dados}).

\textbf{No nosso banco} (previsto pelo edital, ainda não visto na amostra de
provas): granularidade (o grão da fato); \textbf{SCD tipo 2}; as fases do
\textbf{CRISP-DM}; fato semiaditivo; as quatro características de Inmon;
clusterização $\times$ regras de associação; DW $\times$ Data Lake;
drill-down $\times$ roll-up.

Rode \texttt{./quiz.py bi}.
\end{jacaiu}

\section{Alta probabilidade / pesquisa extra}

\begin{itemize}
\item \textbf{Kimball} (bottom-up, data marts) $\times$ \textbf{Inmon}
(top-down, DW corporativo).
\item \textbf{Self-service BI} e \textbf{storytelling com dados} (aparece no
Perfil 1).
\item \textbf{KPI $\times$ métrica:} KPI é indicador-chave atrelado a um
objetivo.
\item Absolutos a descartar: ``ferramentas de BI são incompatíveis com
tabelas normalizadas'', ``3FN maximiza performance de agregação no DW'' ---
falso; 3FN é para OLTP, não para o DW analítico.
\end{itemize}

\begin{comosair}
Decore o par Kimball: Star = 1 fato + dimensões planas (desnormalizadas) =
menos joins = mais rápido para BI. Snowflake = dimensões normalizadas = mais
joins = mais lento, porém economiza espaço. OLAP dice = filtro em MÚLTIPLAS
dimensões; slice = filtro em UMA. Camadas do Lakehouse: Bronze (bruto) $\to$
Silver (limpo) $\to$ Gold (dimensional para BI). Gatilho de erro: ``sempre'',
``obrigatório'', ``incompatível'', ``3FN no DW'' --- desconfie.
\end{comosair}
